% ============================================================
% CHAPITRE : CONTEXTE GÉNÉRAL DU STAGE
% Rapport PFE - Othmane SADIKI
% À inclure dans le document principal avec % ============================================================
% CHAPITRE : CONTEXTE GÉNÉRAL DU STAGE
% Rapport PFE - Othmane SADIKI
% À inclure dans le document principal avec % ============================================================
% CHAPITRE : CONTEXTE GÉNÉRAL DU STAGE
% Rapport PFE - Othmane SADIKI
% À inclure dans le document principal avec % ============================================================
% CHAPITRE : CONTEXTE GÉNÉRAL DU STAGE
% Rapport PFE - Othmane SADIKI
% À inclure dans le document principal avec \input{presentation_entreprise}
% ============================================================

\chapter{Contexte général du stage}

\section{Présentation de l'entreprise d'accueil}

\subsection{Fondation et direction}

Le Groupe Silamir a été fondé en 2011 par Muriel Figer et Juliette Soria, partant du constat que dans un contexte d'immédiateté et d'incertitude, les organisations doivent évoluer en permanence pour rester compétitives. Depuis sa création, Silamir s'est imposé comme un acteur majeur du conseil en transformation digitale, accompagnant les entreprises dans l'optimisation de leurs processus et l'adoption de technologies avancées.

\subsection{Services et domaines d'expertise}

Silamir propose une combinaison intégrée de conseil de bout en bout et de technologies avancées, structurée autour de cinq pôles d'expertise principaux :

\begin{itemize}
    \item \textbf{Intelligence Artificielle et Data \& IA} : exploitation des données et intelligence artificielle au service de la performance métier ;
    \item \textbf{Cybersécurité} : protection des systèmes d'information et gestion des risques numériques ;
    \item \textbf{Excellence Commerciale} (lancée en 2024 sous la marque Silamir Grow) : optimisation de la performance commerciale ;
    \item \textbf{Process Intelligence} : analyse et optimisation des processus métier par le Process Mining, notamment via la plateforme Celonis ;
    \item \textbf{Silamind} (lancé en 2024) : entité dédiée en tant que pure player Microsoft.
\end{itemize}

\subsection{Trajectoire de croissance}

Depuis sa fondation, le Groupe Silamir a connu une croissance soutenue, marquée par plusieurs acquisitions stratégiques et expansions géographiques :

\begin{itemize}
    \item \textbf{2015} : Acquisition de Caiman Design, renforçant les capacités de conception ;
    \item \textbf{2019} : Acquisition de Digital \& You (une ESN) et implantation au Maroc, ouvrant un centre de services nearshore ;
    \item \textbf{2022} : Naxicap Partners acquiert une participation dans le groupe, accélérant la stratégie de croissance ;
    \item \textbf{2023} : Acquisition de Polarys (compétences Data \& IA en France, Pologne et Émirats Arabes Unis) et de FSystem (cybersécurité) ;
    \item \textbf{2024} : Lancement de Silamind (pure player Microsoft) et de Silamir Grow (Excellence Commerciale).
\end{itemize}

% ---- FIGURE : Chronologie de Silamir ----
\begin{figure}[htbp]
    \centering
    % \includegraphics[width=0.9\textwidth]{figures/silamir_timeline.png}
    \fbox{\parbox{0.85\textwidth}{\centering\vspace{2cm}\textit{Figure à insérer : Chronologie de Silamir}\vspace{2cm}}}
    \caption{Chronologie du Groupe Silamir}
    \label{fig:silamir_timeline}
\end{figure}

\subsection{Taille et ambitions}

En 2025, le Groupe Silamir emploie environ 700 collaborateurs et opère dans plus de 30 pays. Le groupe dispose de centres de services nearshore au Maroc et en Pologne. Le siège social est situé à Paris, avec des bureaux régionaux à Casablanca, Cracovie, Nancy, Nantes et Rennes.

Le groupe affiche des ambitions élevées avec un objectif de \textbf{130 millions d'euros de chiffre d'affaires} et un effectif de \textbf{1\,200 collaborateurs} à l'horizon 2026.

\subsection{Reconnaissances et distinctions}

Le Groupe Silamir a été récompensé à plusieurs reprises pour son parcours entrepreneurial et son engagement en faveur de la parité :

\begin{itemize}
    \item \textbf{2021} : Reconnu dans le Top 50 \#WomenEquity ;
    \item \textbf{2022} : Inclus dans le classement French Women Entrepreneurs 40 (\#FWE40) ;
    \item \textbf{2023} : Lauréat du prix \og Start c'est bien, Up c'est mieux \fg{} et la co-fondatrice Juliette Soria choisie pour représenter la France au Sommet des Jeunes Entrepreneurs du G20 ;
    \item \textbf{2024} : Lauréat de la promotion FWE40.
\end{itemize}

\subsection{Siège social et implantations}

Le siège social du groupe est situé à Paris. Les principaux bureaux sont répartis entre Casablanca, Cracovie, Nancy, Nantes et Rennes, offrant une couverture géographique étendue et une proximité avec les clients internationaux.

% ---- FIGURE : Organigramme ----
\begin{figure}[htbp]
    \centering
    % \includegraphics[width=0.85\textwidth]{figures/organigramme.png}
    \fbox{\parbox{0.85\textwidth}{\centering\vspace{3cm}\textit{Figure à insérer : Organigramme du Groupe Silamir}\vspace{3cm}}}
    \caption{Organigramme du Groupe Silamir}
    \label{fig:organigramme}
\end{figure}

% ============================================================
\section{Présentation du département Process Intelligence}
% ============================================================

J'effectue mon stage au sein du département \textbf{Process Intelligence (PI)} de Silamir, un domaine de conseil spécialisé et de pointe qui va au-delà de l'analytique de données traditionnelle. Mon choix a été motivé par le désir de travailler à l'intersection de l'ingénierie des données, du Process Mining et des résultats métier concrets. Ce département offre un chemin direct pour y parvenir en se concentrant sur l'ADN fondamental des entreprises : leurs processus opérationnels.

\subsection{Le Process Intelligence : définition et enjeux}

Le Process Intelligence est une discipline dédiée à l'utilisation de techniques basées sur les données pour \textbf{découvrir, surveiller et optimiser les processus métier}. Il répond non seulement à la question \textit{«~que se passe-t-il~?~»} dans une organisation, mais aussi au \textit{«~comment~»} et au \textit{«~pourquoi~»}. En analysant les empreintes numériques laissées dans les systèmes d'information (ERP, CRM, etc.), le PI fournit une radiographie objective et factuelle de l'exécution des processus, révélant :

\begin{itemize}
    \item Les \textbf{inefficacités} et les \textbf{goulots d'étranglement} dans les flux de travail ;
    \item Les \textbf{écarts de conformité} entre le processus réel et le processus cible ;
    \item Les \textbf{opportunités d'automatisation} souvent invisibles à l'œil nu ;
    \item Les \textbf{causes profondes} des dysfonctionnements opérationnels.
\end{itemize}

\subsection{La plateforme Celonis}

Le travail du département est rendu possible par le \textbf{Celonis Execution Management System (EMS)}, la plateforme leader mondial du Process Mining et de la gestion de l'exécution. Chez Silamir, Celonis n'est pas simplement un outil d'analyse pour consultants ; c'est la \textbf{plateforme centrale} sur laquelle est délivrée la valeur client. La plateforme permet de :

\begin{itemize}
    \item Connecter les données des systèmes source (SAP, Salesforce, Oracle, etc.) ;
    \item Construire des modèles de données (Data Models) représentant les processus métier ;
    \item Découvrir automatiquement les flux de processus réels (Process Discovery) ;
    \item Comparer les processus réels aux processus cibles (Conformance Checking) ;
    \item Créer des KPIs et des tableaux de bord interactifs ;
    \item Déclencher des actions correctives via l'Action Engine.
\end{itemize}

% ---- FIGURE : Architecture Celonis ----
\begin{figure}[htbp]
    \centering
    % \includegraphics[width=0.8\textwidth]{figures/celonis_architecture.png}
    \fbox{\parbox{0.85\textwidth}{\centering\vspace{2.5cm}\textit{Figure à insérer : Architecture de la plateforme Celonis EMS}\vspace{2.5cm}}}
    \caption{Architecture de la plateforme Celonis EMS}
    \label{fig:celonis_architecture}
\end{figure}

\subsection{Composition de l'équipe}

L'équipe dans laquelle je me suis intégré est composée de profils complémentaires, chacun jouant un rôle spécifique dans la chaîne de valeur du Process Intelligence :

\begin{itemize}
    \item \textbf{Data Engineers} : responsables de la construction des pipelines de données, de l'ETL et des modèles de données dans Celonis ;
    \item \textbf{Data Analysts et Consultants Processus} : traduisent les données modélisées en insights métier et identifient les opportunités d'optimisation ;
    \item \textbf{Développeurs RPA} : conçoivent des solutions automatisées pour les inefficacités découvertes par le Process Mining ;
    \item \textbf{Tech Lead} : fournit les orientations architecturales et assure l'alignement avec les objectifs stratégiques du département.
\end{itemize}

\noindent L'équipe projet à laquelle je suis rattaché est constituée de :

\begin{table}[htbp]
    \centering
    \begin{tabular}{|l|l|}
        \hline
        \textbf{Membre} & \textbf{Rôle} \\
        \hline
        Mme Hasna IGHAZRAN & Tech Lead / Data Engineer Senior \\
        \hline
        M. Gabriel [NOM] & Data Engineer / Source des données \\
        \hline
        M. Othmane SADIKI & Stagiaire Data Engineer \& Data Scientist \\
        \hline
    \end{tabular}
    \caption{Composition de l'équipe projet}
    \label{tab:equipe}
\end{table}

% ============================================================
\section{Contexte du projet}
% ============================================================

Mon stage s'inscrit dans le cadre d'un projet de \textbf{Process Intelligence} pour un client utilisant \textbf{SAP} comme ERP. L'objectif est d'analyser et d'optimiser le processus \textbf{Account Payable} (Comptes Fournisseurs) à travers la plateforme Celonis, en combinant Process Mining et Machine Learning prédictif.

\subsection{Problématique métier}

Le client fait face à plusieurs défis liés à son processus de paiement fournisseur :

\begin{itemize}
    \item Des \textbf{retards récurrents} dans le paiement des factures fournisseurs ;
    \item Un \textbf{impact négatif sur le flux de trésorerie} (cash flow) ;
    \item Des \textbf{escomptes non capturés} (cash discount) par manque de réactivité ;
    \item Des risques de \textbf{non-conformité réglementaire} ;
    \item Des \textbf{factures bloquées} dont les causes ne sont pas clairement identifiées ;
    \item Une \textbf{visibilité insuffisante} sur les fournisseurs à risque.
\end{itemize}

\subsection{Objectifs du projet}

L'objectif principal est d'améliorer les délais de paiement fournisseur pour optimiser le flux de trésorerie et assurer la conformité réglementaire. Plus spécifiquement, le projet vise à :

\begin{enumerate}
    \item Extraire, transformer et charger les données SAP dans Celonis (ETL) ;
    \item Construire un Data Model performant du processus Account Payable ;
    \item Créer des KPIs et des tableaux de bord pour analyser les délais de paiement et l'impact cash ;
    \item Réaliser un Conformance Checking entre le processus réel et le processus cible ;
    \item Développer un modèle de Machine Learning pour prédire les retards de paiement ;
    \item Proposer un système d'alertes et de scoring de risque intégrable dans Celonis.
\end{enumerate}

\subsection{Périmètre technique}

Le tableau~\ref{tab:scope_technique} résume le périmètre technique du projet.

\begin{table}[htbp]
    \centering
    \begin{tabular}{|l|p{10cm}|}
        \hline
        \textbf{Élément} & \textbf{Description} \\
        \hline
        ERP source & SAP (modules FI / MM) \\
        \hline
        Tables clés & BSEG (lignes d'écritures comptables), BKPF (en-têtes de pièces comptables), LFA1 (master data fournisseurs), EKKO/EKPO (commandes d'achat), RBKP (factures fournisseurs) \\
        \hline
        Plateforme & Celonis EMS (environnement Sandbox) \\
        \hline
        Langages & Python (pandas, scikit-learn, XGBoost, SHAP), PQL (Process Query Language) \\
        \hline
        Méthodologie & CRISP-DM (Cross-Industry Standard Process for Data Mining) \\
        \hline
    \end{tabular}
    \caption{Périmètre technique du projet}
    \label{tab:scope_technique}
\end{table}

\subsection{KPIs clés prioritaires}

Les indicateurs clés de performance (KPIs) prioritaires identifiés pour ce projet sont :

\begin{itemize}
    \item \textbf{Délai de paiement moyen} : nombre de jours entre la réception de la facture et le paiement effectif ;
    \item \textbf{Taux de retard} : pourcentage de factures payées après la date d'échéance ;
    \item \textbf{Impact Cash} : coût financier des retards et des escomptes non capturés ;
    \item \textbf{Taux d'escomptes capturés} : pourcentage d'escomptes effectivement obtenus vs disponibles ;
    \item \textbf{Taux de factures bloquées} : pourcentage et causes des blocages ;
    \item \textbf{Score de risque fournisseur} : score ML de probabilité de retard (0--100\%).
\end{itemize}

% ---- FIGURE : Processus Account Payable ----
\begin{figure}[htbp]
    \centering
    % \includegraphics[width=0.9\textwidth]{figures/processus_ap.png}
    \fbox{\parbox{0.85\textwidth}{\centering\vspace{3cm}\textit{Figure à insérer : Flux du processus Account Payable cible}\vspace{3cm}}}
    \caption{Flux du processus Account Payable (processus cible)}
    \label{fig:processus_ap}
\end{figure}

% ============================================================

\chapter{Contexte général du stage}

\section{Présentation de l'entreprise d'accueil}

\subsection{Fondation et direction}

Le Groupe Silamir a été fondé en 2011 par Muriel Figer et Juliette Soria, partant du constat que dans un contexte d'immédiateté et d'incertitude, les organisations doivent évoluer en permanence pour rester compétitives. Depuis sa création, Silamir s'est imposé comme un acteur majeur du conseil en transformation digitale, accompagnant les entreprises dans l'optimisation de leurs processus et l'adoption de technologies avancées.

\subsection{Services et domaines d'expertise}

Silamir propose une combinaison intégrée de conseil de bout en bout et de technologies avancées, structurée autour de cinq pôles d'expertise principaux :

\begin{itemize}
    \item \textbf{Intelligence Artificielle et Data \& IA} : exploitation des données et intelligence artificielle au service de la performance métier ;
    \item \textbf{Cybersécurité} : protection des systèmes d'information et gestion des risques numériques ;
    \item \textbf{Excellence Commerciale} (lancée en 2024 sous la marque Silamir Grow) : optimisation de la performance commerciale ;
    \item \textbf{Process Intelligence} : analyse et optimisation des processus métier par le Process Mining, notamment via la plateforme Celonis ;
    \item \textbf{Silamind} (lancé en 2024) : entité dédiée en tant que pure player Microsoft.
\end{itemize}

\subsection{Trajectoire de croissance}

Depuis sa fondation, le Groupe Silamir a connu une croissance soutenue, marquée par plusieurs acquisitions stratégiques et expansions géographiques :

\begin{itemize}
    \item \textbf{2015} : Acquisition de Caiman Design, renforçant les capacités de conception ;
    \item \textbf{2019} : Acquisition de Digital \& You (une ESN) et implantation au Maroc, ouvrant un centre de services nearshore ;
    \item \textbf{2022} : Naxicap Partners acquiert une participation dans le groupe, accélérant la stratégie de croissance ;
    \item \textbf{2023} : Acquisition de Polarys (compétences Data \& IA en France, Pologne et Émirats Arabes Unis) et de FSystem (cybersécurité) ;
    \item \textbf{2024} : Lancement de Silamind (pure player Microsoft) et de Silamir Grow (Excellence Commerciale).
\end{itemize}

% ---- FIGURE : Chronologie de Silamir ----
\begin{figure}[htbp]
    \centering
    % \includegraphics[width=0.9\textwidth]{figures/silamir_timeline.png}
    \fbox{\parbox{0.85\textwidth}{\centering\vspace{2cm}\textit{Figure à insérer : Chronologie de Silamir}\vspace{2cm}}}
    \caption{Chronologie du Groupe Silamir}
    \label{fig:silamir_timeline}
\end{figure}

\subsection{Taille et ambitions}

En 2025, le Groupe Silamir emploie environ 700 collaborateurs et opère dans plus de 30 pays. Le groupe dispose de centres de services nearshore au Maroc et en Pologne. Le siège social est situé à Paris, avec des bureaux régionaux à Casablanca, Cracovie, Nancy, Nantes et Rennes.

Le groupe affiche des ambitions élevées avec un objectif de \textbf{130 millions d'euros de chiffre d'affaires} et un effectif de \textbf{1\,200 collaborateurs} à l'horizon 2026.

\subsection{Reconnaissances et distinctions}

Le Groupe Silamir a été récompensé à plusieurs reprises pour son parcours entrepreneurial et son engagement en faveur de la parité :

\begin{itemize}
    \item \textbf{2021} : Reconnu dans le Top 50 \#WomenEquity ;
    \item \textbf{2022} : Inclus dans le classement French Women Entrepreneurs 40 (\#FWE40) ;
    \item \textbf{2023} : Lauréat du prix \og Start c'est bien, Up c'est mieux \fg{} et la co-fondatrice Juliette Soria choisie pour représenter la France au Sommet des Jeunes Entrepreneurs du G20 ;
    \item \textbf{2024} : Lauréat de la promotion FWE40.
\end{itemize}

\subsection{Siège social et implantations}

Le siège social du groupe est situé à Paris. Les principaux bureaux sont répartis entre Casablanca, Cracovie, Nancy, Nantes et Rennes, offrant une couverture géographique étendue et une proximité avec les clients internationaux.

% ---- FIGURE : Organigramme ----
\begin{figure}[htbp]
    \centering
    % \includegraphics[width=0.85\textwidth]{figures/organigramme.png}
    \fbox{\parbox{0.85\textwidth}{\centering\vspace{3cm}\textit{Figure à insérer : Organigramme du Groupe Silamir}\vspace{3cm}}}
    \caption{Organigramme du Groupe Silamir}
    \label{fig:organigramme}
\end{figure}

% ============================================================
\section{Présentation du département Process Intelligence}
% ============================================================

J'effectue mon stage au sein du département \textbf{Process Intelligence (PI)} de Silamir, un domaine de conseil spécialisé et de pointe qui va au-delà de l'analytique de données traditionnelle. Mon choix a été motivé par le désir de travailler à l'intersection de l'ingénierie des données, du Process Mining et des résultats métier concrets. Ce département offre un chemin direct pour y parvenir en se concentrant sur l'ADN fondamental des entreprises : leurs processus opérationnels.

\subsection{Le Process Intelligence : définition et enjeux}

Le Process Intelligence est une discipline dédiée à l'utilisation de techniques basées sur les données pour \textbf{découvrir, surveiller et optimiser les processus métier}. Il répond non seulement à la question \textit{«~que se passe-t-il~?~»} dans une organisation, mais aussi au \textit{«~comment~»} et au \textit{«~pourquoi~»}. En analysant les empreintes numériques laissées dans les systèmes d'information (ERP, CRM, etc.), le PI fournit une radiographie objective et factuelle de l'exécution des processus, révélant :

\begin{itemize}
    \item Les \textbf{inefficacités} et les \textbf{goulots d'étranglement} dans les flux de travail ;
    \item Les \textbf{écarts de conformité} entre le processus réel et le processus cible ;
    \item Les \textbf{opportunités d'automatisation} souvent invisibles à l'œil nu ;
    \item Les \textbf{causes profondes} des dysfonctionnements opérationnels.
\end{itemize}

\subsection{La plateforme Celonis}

Le travail du département est rendu possible par le \textbf{Celonis Execution Management System (EMS)}, la plateforme leader mondial du Process Mining et de la gestion de l'exécution. Chez Silamir, Celonis n'est pas simplement un outil d'analyse pour consultants ; c'est la \textbf{plateforme centrale} sur laquelle est délivrée la valeur client. La plateforme permet de :

\begin{itemize}
    \item Connecter les données des systèmes source (SAP, Salesforce, Oracle, etc.) ;
    \item Construire des modèles de données (Data Models) représentant les processus métier ;
    \item Découvrir automatiquement les flux de processus réels (Process Discovery) ;
    \item Comparer les processus réels aux processus cibles (Conformance Checking) ;
    \item Créer des KPIs et des tableaux de bord interactifs ;
    \item Déclencher des actions correctives via l'Action Engine.
\end{itemize}

% ---- FIGURE : Architecture Celonis ----
\begin{figure}[htbp]
    \centering
    % \includegraphics[width=0.8\textwidth]{figures/celonis_architecture.png}
    \fbox{\parbox{0.85\textwidth}{\centering\vspace{2.5cm}\textit{Figure à insérer : Architecture de la plateforme Celonis EMS}\vspace{2.5cm}}}
    \caption{Architecture de la plateforme Celonis EMS}
    \label{fig:celonis_architecture}
\end{figure}

\subsection{Composition de l'équipe}

L'équipe dans laquelle je me suis intégré est composée de profils complémentaires, chacun jouant un rôle spécifique dans la chaîne de valeur du Process Intelligence :

\begin{itemize}
    \item \textbf{Data Engineers} : responsables de la construction des pipelines de données, de l'ETL et des modèles de données dans Celonis ;
    \item \textbf{Data Analysts et Consultants Processus} : traduisent les données modélisées en insights métier et identifient les opportunités d'optimisation ;
    \item \textbf{Développeurs RPA} : conçoivent des solutions automatisées pour les inefficacités découvertes par le Process Mining ;
    \item \textbf{Tech Lead} : fournit les orientations architecturales et assure l'alignement avec les objectifs stratégiques du département.
\end{itemize}

\noindent L'équipe projet à laquelle je suis rattaché est constituée de :

\begin{table}[htbp]
    \centering
    \begin{tabular}{|l|l|}
        \hline
        \textbf{Membre} & \textbf{Rôle} \\
        \hline
        Mme Hasna IGHAZRAN & Tech Lead / Data Engineer Senior \\
        \hline
        M. Gabriel [NOM] & Data Engineer / Source des données \\
        \hline
        M. Othmane SADIKI & Stagiaire Data Engineer \& Data Scientist \\
        \hline
    \end{tabular}
    \caption{Composition de l'équipe projet}
    \label{tab:equipe}
\end{table}

% ============================================================
\section{Contexte du projet}
% ============================================================

Mon stage s'inscrit dans le cadre d'un projet de \textbf{Process Intelligence} pour un client utilisant \textbf{SAP} comme ERP. L'objectif est d'analyser et d'optimiser le processus \textbf{Account Payable} (Comptes Fournisseurs) à travers la plateforme Celonis, en combinant Process Mining et Machine Learning prédictif.

\subsection{Problématique métier}

Le client fait face à plusieurs défis liés à son processus de paiement fournisseur :

\begin{itemize}
    \item Des \textbf{retards récurrents} dans le paiement des factures fournisseurs ;
    \item Un \textbf{impact négatif sur le flux de trésorerie} (cash flow) ;
    \item Des \textbf{escomptes non capturés} (cash discount) par manque de réactivité ;
    \item Des risques de \textbf{non-conformité réglementaire} ;
    \item Des \textbf{factures bloquées} dont les causes ne sont pas clairement identifiées ;
    \item Une \textbf{visibilité insuffisante} sur les fournisseurs à risque.
\end{itemize}

\subsection{Objectifs du projet}

L'objectif principal est d'améliorer les délais de paiement fournisseur pour optimiser le flux de trésorerie et assurer la conformité réglementaire. Plus spécifiquement, le projet vise à :

\begin{enumerate}
    \item Extraire, transformer et charger les données SAP dans Celonis (ETL) ;
    \item Construire un Data Model performant du processus Account Payable ;
    \item Créer des KPIs et des tableaux de bord pour analyser les délais de paiement et l'impact cash ;
    \item Réaliser un Conformance Checking entre le processus réel et le processus cible ;
    \item Développer un modèle de Machine Learning pour prédire les retards de paiement ;
    \item Proposer un système d'alertes et de scoring de risque intégrable dans Celonis.
\end{enumerate}

\subsection{Périmètre technique}

Le tableau~\ref{tab:scope_technique} résume le périmètre technique du projet.

\begin{table}[htbp]
    \centering
    \begin{tabular}{|l|p{10cm}|}
        \hline
        \textbf{Élément} & \textbf{Description} \\
        \hline
        ERP source & SAP (modules FI / MM) \\
        \hline
        Tables clés & BSEG (lignes d'écritures comptables), BKPF (en-têtes de pièces comptables), LFA1 (master data fournisseurs), EKKO/EKPO (commandes d'achat), RBKP (factures fournisseurs) \\
        \hline
        Plateforme & Celonis EMS (environnement Sandbox) \\
        \hline
        Langages & Python (pandas, scikit-learn, XGBoost, SHAP), PQL (Process Query Language) \\
        \hline
        Méthodologie & CRISP-DM (Cross-Industry Standard Process for Data Mining) \\
        \hline
    \end{tabular}
    \caption{Périmètre technique du projet}
    \label{tab:scope_technique}
\end{table}

\subsection{KPIs clés prioritaires}

Les indicateurs clés de performance (KPIs) prioritaires identifiés pour ce projet sont :

\begin{itemize}
    \item \textbf{Délai de paiement moyen} : nombre de jours entre la réception de la facture et le paiement effectif ;
    \item \textbf{Taux de retard} : pourcentage de factures payées après la date d'échéance ;
    \item \textbf{Impact Cash} : coût financier des retards et des escomptes non capturés ;
    \item \textbf{Taux d'escomptes capturés} : pourcentage d'escomptes effectivement obtenus vs disponibles ;
    \item \textbf{Taux de factures bloquées} : pourcentage et causes des blocages ;
    \item \textbf{Score de risque fournisseur} : score ML de probabilité de retard (0--100\%).
\end{itemize}

% ---- FIGURE : Processus Account Payable ----
\begin{figure}[htbp]
    \centering
    % \includegraphics[width=0.9\textwidth]{figures/processus_ap.png}
    \fbox{\parbox{0.85\textwidth}{\centering\vspace{3cm}\textit{Figure à insérer : Flux du processus Account Payable cible}\vspace{3cm}}}
    \caption{Flux du processus Account Payable (processus cible)}
    \label{fig:processus_ap}
\end{figure}

% ============================================================

\chapter{Contexte général du stage}

\section{Présentation de l'entreprise d'accueil}

\subsection{Fondation et direction}

Le Groupe Silamir a été fondé en 2011 par Muriel Figer et Juliette Soria, partant du constat que dans un contexte d'immédiateté et d'incertitude, les organisations doivent évoluer en permanence pour rester compétitives. Depuis sa création, Silamir s'est imposé comme un acteur majeur du conseil en transformation digitale, accompagnant les entreprises dans l'optimisation de leurs processus et l'adoption de technologies avancées.

\subsection{Services et domaines d'expertise}

Silamir propose une combinaison intégrée de conseil de bout en bout et de technologies avancées, structurée autour de cinq pôles d'expertise principaux :

\begin{itemize}
    \item \textbf{Intelligence Artificielle et Data \& IA} : exploitation des données et intelligence artificielle au service de la performance métier ;
    \item \textbf{Cybersécurité} : protection des systèmes d'information et gestion des risques numériques ;
    \item \textbf{Excellence Commerciale} (lancée en 2024 sous la marque Silamir Grow) : optimisation de la performance commerciale ;
    \item \textbf{Process Intelligence} : analyse et optimisation des processus métier par le Process Mining, notamment via la plateforme Celonis ;
    \item \textbf{Silamind} (lancé en 2024) : entité dédiée en tant que pure player Microsoft.
\end{itemize}

\subsection{Trajectoire de croissance}

Depuis sa fondation, le Groupe Silamir a connu une croissance soutenue, marquée par plusieurs acquisitions stratégiques et expansions géographiques :

\begin{itemize}
    \item \textbf{2015} : Acquisition de Caiman Design, renforçant les capacités de conception ;
    \item \textbf{2019} : Acquisition de Digital \& You (une ESN) et implantation au Maroc, ouvrant un centre de services nearshore ;
    \item \textbf{2022} : Naxicap Partners acquiert une participation dans le groupe, accélérant la stratégie de croissance ;
    \item \textbf{2023} : Acquisition de Polarys (compétences Data \& IA en France, Pologne et Émirats Arabes Unis) et de FSystem (cybersécurité) ;
    \item \textbf{2024} : Lancement de Silamind (pure player Microsoft) et de Silamir Grow (Excellence Commerciale).
\end{itemize}

% ---- FIGURE : Chronologie de Silamir ----
\begin{figure}[htbp]
    \centering
    % \includegraphics[width=0.9\textwidth]{figures/silamir_timeline.png}
    \fbox{\parbox{0.85\textwidth}{\centering\vspace{2cm}\textit{Figure à insérer : Chronologie de Silamir}\vspace{2cm}}}
    \caption{Chronologie du Groupe Silamir}
    \label{fig:silamir_timeline}
\end{figure}

\subsection{Taille et ambitions}

En 2025, le Groupe Silamir emploie environ 700 collaborateurs et opère dans plus de 30 pays. Le groupe dispose de centres de services nearshore au Maroc et en Pologne. Le siège social est situé à Paris, avec des bureaux régionaux à Casablanca, Cracovie, Nancy, Nantes et Rennes.

Le groupe affiche des ambitions élevées avec un objectif de \textbf{130 millions d'euros de chiffre d'affaires} et un effectif de \textbf{1\,200 collaborateurs} à l'horizon 2026.

\subsection{Reconnaissances et distinctions}

Le Groupe Silamir a été récompensé à plusieurs reprises pour son parcours entrepreneurial et son engagement en faveur de la parité :

\begin{itemize}
    \item \textbf{2021} : Reconnu dans le Top 50 \#WomenEquity ;
    \item \textbf{2022} : Inclus dans le classement French Women Entrepreneurs 40 (\#FWE40) ;
    \item \textbf{2023} : Lauréat du prix \og Start c'est bien, Up c'est mieux \fg{} et la co-fondatrice Juliette Soria choisie pour représenter la France au Sommet des Jeunes Entrepreneurs du G20 ;
    \item \textbf{2024} : Lauréat de la promotion FWE40.
\end{itemize}

\subsection{Siège social et implantations}

Le siège social du groupe est situé à Paris. Les principaux bureaux sont répartis entre Casablanca, Cracovie, Nancy, Nantes et Rennes, offrant une couverture géographique étendue et une proximité avec les clients internationaux.

% ---- FIGURE : Organigramme ----
\begin{figure}[htbp]
    \centering
    % \includegraphics[width=0.85\textwidth]{figures/organigramme.png}
    \fbox{\parbox{0.85\textwidth}{\centering\vspace{3cm}\textit{Figure à insérer : Organigramme du Groupe Silamir}\vspace{3cm}}}
    \caption{Organigramme du Groupe Silamir}
    \label{fig:organigramme}
\end{figure}

% ============================================================
\section{Présentation du département Process Intelligence}
% ============================================================

J'effectue mon stage au sein du département \textbf{Process Intelligence (PI)} de Silamir, un domaine de conseil spécialisé et de pointe qui va au-delà de l'analytique de données traditionnelle. Mon choix a été motivé par le désir de travailler à l'intersection de l'ingénierie des données, du Process Mining et des résultats métier concrets. Ce département offre un chemin direct pour y parvenir en se concentrant sur l'ADN fondamental des entreprises : leurs processus opérationnels.

\subsection{Le Process Intelligence : définition et enjeux}

Le Process Intelligence est une discipline dédiée à l'utilisation de techniques basées sur les données pour \textbf{découvrir, surveiller et optimiser les processus métier}. Il répond non seulement à la question \textit{«~que se passe-t-il~?~»} dans une organisation, mais aussi au \textit{«~comment~»} et au \textit{«~pourquoi~»}. En analysant les empreintes numériques laissées dans les systèmes d'information (ERP, CRM, etc.), le PI fournit une radiographie objective et factuelle de l'exécution des processus, révélant :

\begin{itemize}
    \item Les \textbf{inefficacités} et les \textbf{goulots d'étranglement} dans les flux de travail ;
    \item Les \textbf{écarts de conformité} entre le processus réel et le processus cible ;
    \item Les \textbf{opportunités d'automatisation} souvent invisibles à l'œil nu ;
    \item Les \textbf{causes profondes} des dysfonctionnements opérationnels.
\end{itemize}

\subsection{La plateforme Celonis}

Le travail du département est rendu possible par le \textbf{Celonis Execution Management System (EMS)}, la plateforme leader mondial du Process Mining et de la gestion de l'exécution. Chez Silamir, Celonis n'est pas simplement un outil d'analyse pour consultants ; c'est la \textbf{plateforme centrale} sur laquelle est délivrée la valeur client. La plateforme permet de :

\begin{itemize}
    \item Connecter les données des systèmes source (SAP, Salesforce, Oracle, etc.) ;
    \item Construire des modèles de données (Data Models) représentant les processus métier ;
    \item Découvrir automatiquement les flux de processus réels (Process Discovery) ;
    \item Comparer les processus réels aux processus cibles (Conformance Checking) ;
    \item Créer des KPIs et des tableaux de bord interactifs ;
    \item Déclencher des actions correctives via l'Action Engine.
\end{itemize}

% ---- FIGURE : Architecture Celonis ----
\begin{figure}[htbp]
    \centering
    % \includegraphics[width=0.8\textwidth]{figures/celonis_architecture.png}
    \fbox{\parbox{0.85\textwidth}{\centering\vspace{2.5cm}\textit{Figure à insérer : Architecture de la plateforme Celonis EMS}\vspace{2.5cm}}}
    \caption{Architecture de la plateforme Celonis EMS}
    \label{fig:celonis_architecture}
\end{figure}

\subsection{Composition de l'équipe}

L'équipe dans laquelle je me suis intégré est composée de profils complémentaires, chacun jouant un rôle spécifique dans la chaîne de valeur du Process Intelligence :

\begin{itemize}
    \item \textbf{Data Engineers} : responsables de la construction des pipelines de données, de l'ETL et des modèles de données dans Celonis ;
    \item \textbf{Data Analysts et Consultants Processus} : traduisent les données modélisées en insights métier et identifient les opportunités d'optimisation ;
    \item \textbf{Développeurs RPA} : conçoivent des solutions automatisées pour les inefficacités découvertes par le Process Mining ;
    \item \textbf{Tech Lead} : fournit les orientations architecturales et assure l'alignement avec les objectifs stratégiques du département.
\end{itemize}

\noindent L'équipe projet à laquelle je suis rattaché est constituée de :

\begin{table}[htbp]
    \centering
    \begin{tabular}{|l|l|}
        \hline
        \textbf{Membre} & \textbf{Rôle} \\
        \hline
        Mme Hasna IGHAZRAN & Tech Lead / Data Engineer Senior \\
        \hline
        M. Gabriel [NOM] & Data Engineer / Source des données \\
        \hline
        M. Othmane SADIKI & Stagiaire Data Engineer \& Data Scientist \\
        \hline
    \end{tabular}
    \caption{Composition de l'équipe projet}
    \label{tab:equipe}
\end{table}

% ============================================================
\section{Contexte du projet}
% ============================================================

Mon stage s'inscrit dans le cadre d'un projet de \textbf{Process Intelligence} pour un client utilisant \textbf{SAP} comme ERP. L'objectif est d'analyser et d'optimiser le processus \textbf{Account Payable} (Comptes Fournisseurs) à travers la plateforme Celonis, en combinant Process Mining et Machine Learning prédictif.

\subsection{Problématique métier}

Le client fait face à plusieurs défis liés à son processus de paiement fournisseur :

\begin{itemize}
    \item Des \textbf{retards récurrents} dans le paiement des factures fournisseurs ;
    \item Un \textbf{impact négatif sur le flux de trésorerie} (cash flow) ;
    \item Des \textbf{escomptes non capturés} (cash discount) par manque de réactivité ;
    \item Des risques de \textbf{non-conformité réglementaire} ;
    \item Des \textbf{factures bloquées} dont les causes ne sont pas clairement identifiées ;
    \item Une \textbf{visibilité insuffisante} sur les fournisseurs à risque.
\end{itemize}

\subsection{Objectifs du projet}

L'objectif principal est d'améliorer les délais de paiement fournisseur pour optimiser le flux de trésorerie et assurer la conformité réglementaire. Plus spécifiquement, le projet vise à :

\begin{enumerate}
    \item Extraire, transformer et charger les données SAP dans Celonis (ETL) ;
    \item Construire un Data Model performant du processus Account Payable ;
    \item Créer des KPIs et des tableaux de bord pour analyser les délais de paiement et l'impact cash ;
    \item Réaliser un Conformance Checking entre le processus réel et le processus cible ;
    \item Développer un modèle de Machine Learning pour prédire les retards de paiement ;
    \item Proposer un système d'alertes et de scoring de risque intégrable dans Celonis.
\end{enumerate}

\subsection{Périmètre technique}

Le tableau~\ref{tab:scope_technique} résume le périmètre technique du projet.

\begin{table}[htbp]
    \centering
    \begin{tabular}{|l|p{10cm}|}
        \hline
        \textbf{Élément} & \textbf{Description} \\
        \hline
        ERP source & SAP (modules FI / MM) \\
        \hline
        Tables clés & BSEG (lignes d'écritures comptables), BKPF (en-têtes de pièces comptables), LFA1 (master data fournisseurs), EKKO/EKPO (commandes d'achat), RBKP (factures fournisseurs) \\
        \hline
        Plateforme & Celonis EMS (environnement Sandbox) \\
        \hline
        Langages & Python (pandas, scikit-learn, XGBoost, SHAP), PQL (Process Query Language) \\
        \hline
        Méthodologie & CRISP-DM (Cross-Industry Standard Process for Data Mining) \\
        \hline
    \end{tabular}
    \caption{Périmètre technique du projet}
    \label{tab:scope_technique}
\end{table}

\subsection{KPIs clés prioritaires}

Les indicateurs clés de performance (KPIs) prioritaires identifiés pour ce projet sont :

\begin{itemize}
    \item \textbf{Délai de paiement moyen} : nombre de jours entre la réception de la facture et le paiement effectif ;
    \item \textbf{Taux de retard} : pourcentage de factures payées après la date d'échéance ;
    \item \textbf{Impact Cash} : coût financier des retards et des escomptes non capturés ;
    \item \textbf{Taux d'escomptes capturés} : pourcentage d'escomptes effectivement obtenus vs disponibles ;
    \item \textbf{Taux de factures bloquées} : pourcentage et causes des blocages ;
    \item \textbf{Score de risque fournisseur} : score ML de probabilité de retard (0--100\%).
\end{itemize}

% ---- FIGURE : Processus Account Payable ----
\begin{figure}[htbp]
    \centering
    % \includegraphics[width=0.9\textwidth]{figures/processus_ap.png}
    \fbox{\parbox{0.85\textwidth}{\centering\vspace{3cm}\textit{Figure à insérer : Flux du processus Account Payable cible}\vspace{3cm}}}
    \caption{Flux du processus Account Payable (processus cible)}
    \label{fig:processus_ap}
\end{figure}

% ============================================================

\chapter{Contexte général du stage}

\section{Présentation de l'entreprise d'accueil}

\subsection{Fondation et direction}

Le Groupe Silamir a été fondé en 2011 par Muriel Figer et Juliette Soria, partant du constat que dans un contexte d'immédiateté et d'incertitude, les organisations doivent évoluer en permanence pour rester compétitives. Depuis sa création, Silamir s'est imposé comme un acteur majeur du conseil en transformation digitale, accompagnant les entreprises dans l'optimisation de leurs processus et l'adoption de technologies avancées.

\subsection{Services et domaines d'expertise}

Silamir propose une combinaison intégrée de conseil de bout en bout et de technologies avancées, structurée autour de cinq pôles d'expertise principaux :

\begin{itemize}
    \item \textbf{Intelligence Artificielle et Data \& IA} : exploitation des données et intelligence artificielle au service de la performance métier ;
    \item \textbf{Cybersécurité} : protection des systèmes d'information et gestion des risques numériques ;
    \item \textbf{Excellence Commerciale} (lancée en 2024 sous la marque Silamir Grow) : optimisation de la performance commerciale ;
    \item \textbf{Process Intelligence} : analyse et optimisation des processus métier par le Process Mining, notamment via la plateforme Celonis ;
    \item \textbf{Silamind} (lancé en 2024) : entité dédiée en tant que pure player Microsoft.
\end{itemize}

\subsection{Trajectoire de croissance}

Depuis sa fondation, le Groupe Silamir a connu une croissance soutenue, marquée par plusieurs acquisitions stratégiques et expansions géographiques :

\begin{itemize}
    \item \textbf{2015} : Acquisition de Caiman Design, renforçant les capacités de conception ;
    \item \textbf{2019} : Acquisition de Digital \& You (une ESN) et implantation au Maroc, ouvrant un centre de services nearshore ;
    \item \textbf{2022} : Naxicap Partners acquiert une participation dans le groupe, accélérant la stratégie de croissance ;
    \item \textbf{2023} : Acquisition de Polarys (compétences Data \& IA en France, Pologne et Émirats Arabes Unis) et de FSystem (cybersécurité) ;
    \item \textbf{2024} : Lancement de Silamind (pure player Microsoft) et de Silamir Grow (Excellence Commerciale).
\end{itemize}

% ---- FIGURE : Chronologie de Silamir ----
\begin{figure}[htbp]
    \centering
    % \includegraphics[width=0.9\textwidth]{figures/silamir_timeline.png}
    \fbox{\parbox{0.85\textwidth}{\centering\vspace{2cm}\textit{Figure à insérer : Chronologie de Silamir}\vspace{2cm}}}
    \caption{Chronologie du Groupe Silamir}
    \label{fig:silamir_timeline}
\end{figure}

\subsection{Taille et ambitions}

En 2025, le Groupe Silamir emploie environ 700 collaborateurs et opère dans plus de 30 pays. Le groupe dispose de centres de services nearshore au Maroc et en Pologne. Le siège social est situé à Paris, avec des bureaux régionaux à Casablanca, Cracovie, Nancy, Nantes et Rennes.

Le groupe affiche des ambitions élevées avec un objectif de \textbf{130 millions d'euros de chiffre d'affaires} et un effectif de \textbf{1\,200 collaborateurs} à l'horizon 2026.

\subsection{Reconnaissances et distinctions}

Le Groupe Silamir a été récompensé à plusieurs reprises pour son parcours entrepreneurial et son engagement en faveur de la parité :

\begin{itemize}
    \item \textbf{2021} : Reconnu dans le Top 50 \#WomenEquity ;
    \item \textbf{2022} : Inclus dans le classement French Women Entrepreneurs 40 (\#FWE40) ;
    \item \textbf{2023} : Lauréat du prix \og Start c'est bien, Up c'est mieux \fg{} et la co-fondatrice Juliette Soria choisie pour représenter la France au Sommet des Jeunes Entrepreneurs du G20 ;
    \item \textbf{2024} : Lauréat de la promotion FWE40.
\end{itemize}

\subsection{Siège social et implantations}

Le siège social du groupe est situé à Paris. Les principaux bureaux sont répartis entre Casablanca, Cracovie, Nancy, Nantes et Rennes, offrant une couverture géographique étendue et une proximité avec les clients internationaux.

% ---- FIGURE : Organigramme ----
\begin{figure}[htbp]
    \centering
    % \includegraphics[width=0.85\textwidth]{figures/organigramme.png}
    \fbox{\parbox{0.85\textwidth}{\centering\vspace{3cm}\textit{Figure à insérer : Organigramme du Groupe Silamir}\vspace{3cm}}}
    \caption{Organigramme du Groupe Silamir}
    \label{fig:organigramme}
\end{figure}

% ============================================================
\section{Présentation du département Process Intelligence}
% ============================================================

J'effectue mon stage au sein du département \textbf{Process Intelligence (PI)} de Silamir, un domaine de conseil spécialisé et de pointe qui va au-delà de l'analytique de données traditionnelle. Mon choix a été motivé par le désir de travailler à l'intersection de l'ingénierie des données, du Process Mining et des résultats métier concrets. Ce département offre un chemin direct pour y parvenir en se concentrant sur l'ADN fondamental des entreprises : leurs processus opérationnels.

\subsection{Le Process Intelligence : définition et enjeux}

Le Process Intelligence est une discipline dédiée à l'utilisation de techniques basées sur les données pour \textbf{découvrir, surveiller et optimiser les processus métier}. Il répond non seulement à la question \textit{«~que se passe-t-il~?~»} dans une organisation, mais aussi au \textit{«~comment~»} et au \textit{«~pourquoi~»}. En analysant les empreintes numériques laissées dans les systèmes d'information (ERP, CRM, etc.), le PI fournit une radiographie objective et factuelle de l'exécution des processus, révélant :

\begin{itemize}
    \item Les \textbf{inefficacités} et les \textbf{goulots d'étranglement} dans les flux de travail ;
    \item Les \textbf{écarts de conformité} entre le processus réel et le processus cible ;
    \item Les \textbf{opportunités d'automatisation} souvent invisibles à l'œil nu ;
    \item Les \textbf{causes profondes} des dysfonctionnements opérationnels.
\end{itemize}

\subsection{La plateforme Celonis}

Le travail du département est rendu possible par le \textbf{Celonis Execution Management System (EMS)}, la plateforme leader mondial du Process Mining et de la gestion de l'exécution. Chez Silamir, Celonis n'est pas simplement un outil d'analyse pour consultants ; c'est la \textbf{plateforme centrale} sur laquelle est délivrée la valeur client. La plateforme permet de :

\begin{itemize}
    \item Connecter les données des systèmes source (SAP, Salesforce, Oracle, etc.) ;
    \item Construire des modèles de données (Data Models) représentant les processus métier ;
    \item Découvrir automatiquement les flux de processus réels (Process Discovery) ;
    \item Comparer les processus réels aux processus cibles (Conformance Checking) ;
    \item Créer des KPIs et des tableaux de bord interactifs ;
    \item Déclencher des actions correctives via l'Action Engine.
\end{itemize}

% ---- FIGURE : Architecture Celonis ----
\begin{figure}[htbp]
    \centering
    % \includegraphics[width=0.8\textwidth]{figures/celonis_architecture.png}
    \fbox{\parbox{0.85\textwidth}{\centering\vspace{2.5cm}\textit{Figure à insérer : Architecture de la plateforme Celonis EMS}\vspace{2.5cm}}}
    \caption{Architecture de la plateforme Celonis EMS}
    \label{fig:celonis_architecture}
\end{figure}

\subsection{Composition de l'équipe}

L'équipe dans laquelle je me suis intégré est composée de profils complémentaires, chacun jouant un rôle spécifique dans la chaîne de valeur du Process Intelligence :

\begin{itemize}
    \item \textbf{Data Engineers} : responsables de la construction des pipelines de données, de l'ETL et des modèles de données dans Celonis ;
    \item \textbf{Data Analysts et Consultants Processus} : traduisent les données modélisées en insights métier et identifient les opportunités d'optimisation ;
    \item \textbf{Développeurs RPA} : conçoivent des solutions automatisées pour les inefficacités découvertes par le Process Mining ;
    \item \textbf{Tech Lead} : fournit les orientations architecturales et assure l'alignement avec les objectifs stratégiques du département.
\end{itemize}

\noindent L'équipe projet à laquelle je suis rattaché est constituée de :

\begin{table}[htbp]
    \centering
    \begin{tabular}{|l|l|}
        \hline
        \textbf{Membre} & \textbf{Rôle} \\
        \hline
        Mme Hasna IGHAZRAN & Tech Lead / Data Engineer Senior \\
        \hline
        M. Gabriel [NOM] & Data Engineer / Source des données \\
        \hline
        M. Othmane SADIKI & Stagiaire Data Engineer \& Data Scientist \\
        \hline
    \end{tabular}
    \caption{Composition de l'équipe projet}
    \label{tab:equipe}
\end{table}

% ============================================================
\section{Contexte du projet}
% ============================================================

Mon stage s'inscrit dans le cadre d'un projet de \textbf{Process Intelligence} pour un client utilisant \textbf{SAP} comme ERP. L'objectif est d'analyser et d'optimiser le processus \textbf{Account Payable} (Comptes Fournisseurs) à travers la plateforme Celonis, en combinant Process Mining et Machine Learning prédictif.

\subsection{Problématique métier}

Le client fait face à plusieurs défis liés à son processus de paiement fournisseur :

\begin{itemize}
    \item Des \textbf{retards récurrents} dans le paiement des factures fournisseurs ;
    \item Un \textbf{impact négatif sur le flux de trésorerie} (cash flow) ;
    \item Des \textbf{escomptes non capturés} (cash discount) par manque de réactivité ;
    \item Des risques de \textbf{non-conformité réglementaire} ;
    \item Des \textbf{factures bloquées} dont les causes ne sont pas clairement identifiées ;
    \item Une \textbf{visibilité insuffisante} sur les fournisseurs à risque.
\end{itemize}

\subsection{Objectifs du projet}

L'objectif principal est d'améliorer les délais de paiement fournisseur pour optimiser le flux de trésorerie et assurer la conformité réglementaire. Plus spécifiquement, le projet vise à :

\begin{enumerate}
    \item Extraire, transformer et charger les données SAP dans Celonis (ETL) ;
    \item Construire un Data Model performant du processus Account Payable ;
    \item Créer des KPIs et des tableaux de bord pour analyser les délais de paiement et l'impact cash ;
    \item Réaliser un Conformance Checking entre le processus réel et le processus cible ;
    \item Développer un modèle de Machine Learning pour prédire les retards de paiement ;
    \item Proposer un système d'alertes et de scoring de risque intégrable dans Celonis.
\end{enumerate}

\subsection{Périmètre technique}

Le tableau~\ref{tab:scope_technique} résume le périmètre technique du projet.

\begin{table}[htbp]
    \centering
    \begin{tabular}{|l|p{10cm}|}
        \hline
        \textbf{Élément} & \textbf{Description} \\
        \hline
        ERP source & SAP (modules FI / MM) \\
        \hline
        Tables clés & BSEG (lignes d'écritures comptables), BKPF (en-têtes de pièces comptables), LFA1 (master data fournisseurs), EKKO/EKPO (commandes d'achat), RBKP (factures fournisseurs) \\
        \hline
        Plateforme & Celonis EMS (environnement Sandbox) \\
        \hline
        Langages & Python (pandas, scikit-learn, XGBoost, SHAP), PQL (Process Query Language) \\
        \hline
        Méthodologie & CRISP-DM (Cross-Industry Standard Process for Data Mining) \\
        \hline
    \end{tabular}
    \caption{Périmètre technique du projet}
    \label{tab:scope_technique}
\end{table}

\subsection{KPIs clés prioritaires}

Les indicateurs clés de performance (KPIs) prioritaires identifiés pour ce projet sont :

\begin{itemize}
    \item \textbf{Délai de paiement moyen} : nombre de jours entre la réception de la facture et le paiement effectif ;
    \item \textbf{Taux de retard} : pourcentage de factures payées après la date d'échéance ;
    \item \textbf{Impact Cash} : coût financier des retards et des escomptes non capturés ;
    \item \textbf{Taux d'escomptes capturés} : pourcentage d'escomptes effectivement obtenus vs disponibles ;
    \item \textbf{Taux de factures bloquées} : pourcentage et causes des blocages ;
    \item \textbf{Score de risque fournisseur} : score ML de probabilité de retard (0--100\%).
\end{itemize}

% ---- FIGURE : Processus Account Payable ----
\begin{figure}[htbp]
    \centering
    % \includegraphics[width=0.9\textwidth]{figures/processus_ap.png}
    \fbox{\parbox{0.85\textwidth}{\centering\vspace{3cm}\textit{Figure à insérer : Flux du processus Account Payable cible}\vspace{3cm}}}
    \caption{Flux du processus Account Payable (processus cible)}
    \label{fig:processus_ap}
\end{figure}
